\section{Source history and retained comparison routes}\label{sec:history}
This edition integrates the canonical v0.14 analytic result with its prerequisite arithmetic and transport proofs. The complete original history is retained under \texttt{sources/}; this appendix records its roles without promoting an abandoned route or an old empirical status to a current theorem.

\begin{longtable}{p{.14\linewidth}p{.43\linewidth}p{.34\linewidth}}
\toprule Source stage & Mathematical role retained & Boundary or correction\\\midrule\endhead
Early CM/Gauss notes & Fixed field, Gauss branch, unit presentation and seven-torsion CM comparison & No gamma phase identity obtained by recognition\\
Stark modulus & Conductor $3J$, character normalization, Epstein derivative and unit norm comparison & Modulus equality does not determine complex phase\\
v0.1 & Multiplication norms and the integral doubling projector & Its deleted nonzero component cannot be recovered by repeated symmetric norms\\
v0.2 & Signed eight-point unit packet and explicit carries & Higher-order channels do not remove the first-order phase normalization\\
v0.3 & Exact fifth-torsion discriminant normalizer & Rational elliptic coefficients at the cubic period are still analytic\\
v0.4 & Fifth-isogeny differential and quotient map & Interchart homothety is not a finite CM isogeny identification\\
v0.5 & Conductor Green endpoint and phase-log lift & Conditional branch selection requires an exact phase-sensitive comparison\\
v0.6 & Four-subfield norms and algebraic polarization & The remaining analytic theorem can be stated as exact reality\\
v0.7 & Oriented seven-torsion fixed sectors of a CM threefold & A common threefold does not automatically descend a scalar gamma section\\
v0.8 & Hodge obstruction, complementary torsion and Gauss character & A unit-invariant ordinary line bundle class is not supplied by flat labels alone\\
v0.9 & Explicit seven-dimensional projective descent packet & It does not construct the required meromorphic gamma section\\
v0.10 & Gauge-invariant class test and based conductor holonomy & Raw mixed commutators on the Stein cover are not holomorphic class invariants\\
v0.11 & Fourier orbit obstruction and conjugation-only criterion & No finite periodic-point test certifies global flattening\\
v0.12 & Actual seven-product and second-action transfer rigidity & The raw gamma principal norm is false; its correction matters\\
v0.13 & Full ceiling divisor, transitions and three-theta anchor & Naive conjugation still has unmatched meromorphic orders\\
v0.14 & Cubic growth, genus-three regularization and special cube & Plane regularity and arithmetic evaluation remain distinct\\\bottomrule
\end{longtable}

\subsection{The scalar missing from a closed finite orbit}
The earlier index-23 and index-10 orbit decompositions are $529=1+528$ and $100=1+3+8\cdot12$. They are exact finite permutation data. For values $z_j$ on a cycle of length $d$, replacing every value by $cz_j$ leaves all adjacent ratios unchanged, including their product one, but multiplies the cycle product by $c^d$. An anchor on another fixed-point orbit does not remove this freedom. Thus orbit closure alone does not evaluate a primitive analytic sector or fix its scalar. The modulus proof instead uses Artin induction, the regulator and normalized CM Green products; the complex phase still requires an additional comparison.

\subsection{Three geometries must stay distinct}
The scalar plane carries the meromorphic correction and its canonical divisor product. The intermediate cover $X=\mathbb C^3/\mathbb Z^3$ carries periodic multipliers and a degree-seven principal cover. The compact CM threefold has lattice $\mathbb Z[\omega]^3$ and degree 49 for the same matrix. Finite projective representations or oriented torsion sectors can encode cocycle data in the compact model. They do not alone provide a scalar meromorphic section with the divisor of $R$, a periodic relative multiplier or a based arithmetic identification. The final regularizer resolves the first of these scalar-plane requirements while leaving the others explicit.

\subsection{Editable source and provenance}
The manuscript uses source anchors N (theta normalization), T (norm transport), A (arithmetic modulus), P (principal divisor), D (canonical reconstruction), Q (polarization) and S (conditional rigidity). \texttt{CLAIM\_SOURCE\_INVENTORY.json} records original hashes and copied span counts; the scoped proofread report records the manual mathematical transcription and audit status. Four v0.2 replay-receipt output hashes refer to payloads that were not preserved and remain unverified; intact archived manifests do not repair those missing receipts. Fresh v0.14 evidence is independent. Every diagram caption distinguishes an actual finite family from a schematic or analytic bound. Original verification programs and outputs retain their provenance under their source directories; the verification companions report fresh runs separately.
